% GENERATED FILE. Edit canonical lessons, then run npm run generate:books.
% canonical-source-hash: fnv1a64:8ac9b901f8ba1111
% canonical-lessons: MR-C180-lekhak, MR-C180-citrakar, MR-C180-chayacitrakar, MR-C180-patrakar, MR-C180-vastuvisarad

\chapter{Writer, Painter, Photographer, Journalist, Architect}
\label{ch:mr-writer-painter-photographer-journalist-architect}
\hlchaptermodality{180}

\begin{chapteropening}
\textbf{By the end of this chapter:} \emph{I can say a writer, a painter, a photographer, a journalist and an architect.}

Complete the last lesson of chapter 180: I can say a writer, a painter, a photographer, a journalist and an architect.
\end{chapteropening}

\section[\texorpdfstring{\mr{लेखक} / \mr{लेखिका} (lekhak / lekhikā)}{लेखक / लेखिका (lekhak / lekhikā)}]{\mr{लेखक} / \mr{लेखिका} (lekhak / lekhikā) — a writer}

\label{lesson:MR-C180-lekhak}



\begin{quote}
Before the new one: say the Marathi for an election, then the Marathi for politics.
\end{quote}

\subsection*{You'll want to know: \mr{लेखक} / \mr{लेखिका}}
\textbf{\mr{लेखक} / \mr{लेखिका}} — \emph{lekhak / lekhikā} — \textquotedblleft{}a writer\textquotedblright{}.

\begin{grammarlens}[title={What you've built}]
One more word for this chapter.
\end{grammarlens}

\subsection*{Your turn}
\begin{itemize}
\raggedright
  \item \emph{Say it:} \emph{lekhak / lekhikā}
  \item \emph{Say it:} \emph{lekhak / lekhikā}, once more
  \item \emph{Say it:} say \emph{lekhak / lekhikā}
  \item \emph{Recall:} say the Marathi for an election, then the Marathi for politics, then say \emph{lekhak / lekhikā} again
\end{itemize}

\begin{tcolorbox}[breakable,colback=teal!4,colframe=teal!35!black,title={Before you move on}]
What does \mr{लेखक} / \mr{लेखिका} mean? (\textquotedblleft{}A writer\textquotedblright{}.) Say it once more.
\end{tcolorbox}
\section[\texorpdfstring{\mr{चित्रकार} (citrakār)}{चित्रकार (citrakār)}]{\mr{चित्रकार} (citrakār) — a painter}

\label{lesson:MR-C180-citrakar}



\begin{quote}
Before the new one: say the Marathi for politics, then the Marathi for a writer.
\end{quote}

\subsection*{You'll want to know: \mr{चित्रकार}}
\textbf{\mr{चित्रकार}} — \emph{citrakār} — \textquotedblleft{}a painter\textquotedblright{}.

\begin{grammarlens}[title={What you've built}]
One more word for this chapter.
\end{grammarlens}

\subsection*{Your turn}
\begin{itemize}
\raggedright
  \item \emph{Say it:} \emph{citrakār}
  \item \emph{Say it:} \emph{citrakār}, once more
  \item \emph{Say it:} say \emph{citrakār}
  \item \emph{Recall:} say the Marathi for politics, then the Marathi for a writer, then say \emph{citrakār} again
\end{itemize}

\begin{tcolorbox}[breakable,colback=teal!4,colframe=teal!35!black,title={Before you move on}]
What does \mr{चित्रकार} mean? (\textquotedblleft{}A painter\textquotedblright{}.) Say it once more.
\end{tcolorbox}
\section[\texorpdfstring{\mr{छायाचित्रकार} (chāyācitrakār)}{छायाचित्रकार (chāyācitrakār)}]{\mr{छायाचित्रकार} (chāyācitrakār) — a photographer}

\label{lesson:MR-C180-chayacitrakar}



\begin{quote}
Before the new one: say the Marathi for a writer, then the Marathi for a painter.
\end{quote}

\subsection*{You'll want to know: \mr{छायाचित्रकार}}
\textbf{\mr{छायाचित्रकार}} — \emph{chāyācitrakār} — \textquotedblleft{}a photographer\textquotedblright{}.

\begin{grammarlens}[title={What you've built}]
One more word for this chapter.
\end{grammarlens}

\subsection*{Your turn}
\begin{itemize}
\raggedright
  \item \emph{Say it:} \emph{chāyācitrakār}
  \item \emph{Say it:} \emph{chāyācitrakār}, once more
  \item \emph{Say it:} say \emph{chāyācitrakār}
  \item \emph{Recall:} say the Marathi for a writer, then the Marathi for a painter, then say \emph{chāyācitrakār} again
\end{itemize}

\begin{tcolorbox}[breakable,colback=teal!4,colframe=teal!35!black,title={Before you move on}]
What does \mr{छायाचित्रकार} mean? (\textquotedblleft{}A photographer\textquotedblright{}.) Say it once more.
\end{tcolorbox}
\section[\texorpdfstring{\mr{पत्रकार} (patrakār)}{पत्रकार (patrakār)}]{\mr{पत्रकार} (patrakār) — a journalist}

\label{lesson:MR-C180-patrakar}



\begin{quote}
Before the new one: say the Marathi for a painter, then the Marathi for a photographer.
\end{quote}

\subsection*{You'll want to know: \mr{पत्रकार}}
\textbf{\mr{पत्रकार}} — \emph{patrakār} — \textquotedblleft{}a journalist\textquotedblright{}.

\begin{grammarlens}[title={What you've built}]
One more word for this chapter.
\end{grammarlens}

\subsection*{Your turn}
\begin{itemize}
\raggedright
  \item \emph{Say it:} \emph{patrakār}
  \item \emph{Say it:} \emph{patrakār}, once more
  \item \emph{Say it:} say \emph{patrakār}
  \item \emph{Recall:} say the Marathi for a painter, then the Marathi for a photographer, then say \emph{patrakār} again
\end{itemize}

\begin{tcolorbox}[breakable,colback=teal!4,colframe=teal!35!black,title={Before you move on}]
What does \mr{पत्रकार} mean? (\textquotedblleft{}A journalist\textquotedblright{}.) Say it once more.
\end{tcolorbox}
\section[\texorpdfstring{\mr{वास्तुविशारद} (vāstuviśārad)}{वास्तुविशारद (vāstuviśārad)}]{\mr{वास्तुविशारद} (vāstuviśārad) — an architect}

\label{lesson:MR-C180-vastuvisarad}



\begin{quote}
Before the new one: say the Marathi for a photographer, then the Marathi for a journalist.
\end{quote}

\subsection*{You'll want to know: \mr{वास्तुविशारद}}
\textbf{\mr{वास्तुविशारद}} — \emph{vāstuviśārad} — \textquotedblleft{}an architect\textquotedblright{}.

\begin{grammarlens}[title={What you've built}]
One more word for this chapter.
\end{grammarlens}

\subsection*{Your turn}
\begin{itemize}
\raggedright
  \item \emph{Say it:} \emph{vāstuviśārad}
  \item \emph{Say it:} \emph{vāstuviśārad}, once more
  \item \emph{Say it:} say \emph{vāstuviśārad}
  \item \emph{Recall:} say the Marathi for a photographer, then the Marathi for a journalist, then say \emph{vāstuviśārad} again
\end{itemize}

\begin{tcolorbox}[breakable,colback=teal!4,colframe=teal!35!black,title={Before you move on}]
What does \mr{वास्तुविशारद} mean? (\textquotedblleft{}An architect\textquotedblright{}.) Say it once more.
\end{tcolorbox}
